\ifdefined\pdfobjcompresslevel
\pdfobjcompresslevel=0
\fi
\documentclass[11pt]{article}
\usepackage[margin=1in]{geometry}
\usepackage{amsmath}
\usepackage{amssymb}
\usepackage{graphicx}
\usepackage{tikz-cd}
\usepackage{multicol}
\usepackage{url}
\usepackage{booktabs}
\usepackage{multirow}
\usepackage{makecell}

\setlength{\parindent}{0pt}
\setlength{\parskip}{1\baselineskip}

\begin{document}

\section*{Introduction}

Cerebrovascular disease represents one of the leading causes of death and long-term disability worldwide. In 2021, stroke alone accounted for an estimated 93.8 million prevalent cases and 7.3 million deaths, constituting over 10\% of global mortality \cite{wso2025,gbd2021stroke}. Beyond acute cerebrovascular events, a spectrum of chronic intracranial pathologies --- including unruptured intracranial aneurysms, arteriovenous malformations, and stenoocclusive disease --- carries substantial risk of catastrophic neurological outcomes. A critical, and still incompletely resolved, challenge in the management of these conditions is the reliable, non-invasive quantification of intracranial hemodynamics. Blood flow velocity, wall shear stress (WSS), and pressure distributions within the cerebral vasculature are increasingly recognised as key biomechanical determinants of lesion formation, growth, and rupture risk, yet they remain inaccessible to conventional clinical imaging \cite{wahlin2022,cebral2019wss}.

Four-dimensional Flow MRI (4D Flow MRI) has emerged as the principal non-invasive modality for comprehensive, time-resolved haemodynamic assessment of the cerebral circulation \cite{markl2012,dyverfeldt2015}. By encoding three-directional velocity information across the entire cerebrovascular tree within a single acquisition, 4D Flow MRI enables simultaneous evaluation of arterial and venous flow, net flow quantification, retrospective vessel-of-interest analysis, and the derivation of advanced haemodynamic metrics including WSS and oscillatory shear index \cite{bock2010,stienen2020wss}. Consequently, 4D Flow MRI is increasingly applied to a broad range of cerebrovascular conditions, from the haemodynamic characterisation of intracranial aneurysms and carotid bifurcation disease to the investigation of cerebral autoregulation impairment and flow alterations in neurodegenerative conditions \cite{wahlin2022,sun2020cfd}.

Despite its clinical promise, the practical utility of 4D Flow MRI at standard clinical field strengths (1.5T and 3T) is significantly constrained by its inherently limited signal-to-noise ratio (SNR), spatial resolution, and velocity-to-noise ratio (VNR). These limitations are particularly pronounced in the intracranial circulation, where vessel diameters are small, flow velocities are comparatively low, and accurate haemodynamic quantification demands sub-millimetre spatial resolution \cite{markl2012,dyverfeldt2015}. The partial-volume effect introduced by low spatial resolution systematically biases velocity and WSS estimates, underestimating peak flow velocities and distorting the spatial topology of the velocity field \cite{sun2020cfd}. Furthermore, the extended acquisition times required for whole-brain 4D Flow coverage at 3T compound physiological noise contributions and limit clinical feasibility in acutely ill patients.

Ultra-high-field MRI at 7 Tesla (7T) addresses many of these fundamental limitations. The approximately linear increase in SNR with field strength, combined with improved susceptibility contrast, yields substantially enhanced spatial resolution and vascular conspicuity compared to 3T \cite{kara2022_7t,smra2024}. In the context of 4D Flow MRI, these gains translate directly into higher VNR, reduced partial-volume contamination, and the ability to resolve smaller intracranial vessels --- enabling, for the first time, haemodynamic characterisation at the level of perforator arteries and the distal cortical circulation. Recent work has demonstrated that 7T 4D Flow MRI can achieve repeatable, quantitatively accurate intracranial flow measurements within clinically acceptable acquisition times \cite{mccormick2020_7t4dflow}. However, the limited global availability of 7T systems, the associated infrastructure and safety requirements, and the challenges of integrating investigational ultra-high-field acquisitions into routine clinical workflows place severe constraints on the direct clinical translation of 7T 4D Flow MRI.

This gap between the quality achievable at 7T and the practicality of routine 3T acquisition motivates a computational approach: training deep neural networks to learn the mapping from 3T 4D Flow MRI data to their 7T counterparts, thereby enabling virtual 7T quality from widely available clinical hardware. Deep learning has demonstrated remarkable success in MRI super-resolution and enhancement tasks, learning to recover fine structural detail from low-resolution or low-SNR inputs \cite{duffy2022sr,wang2024srrt}. However, existing super-resolution approaches for 4D Flow MRI are typically trained and validated on synthetic downsampling pipelines that do not capture the complex inter-scanner differences arising from distinct hardware, reconstruction algorithms, and k-space sampling patterns between field strengths \cite{ferdian2020,elahmar2023}. We address this by formulating two complementary learning tasks: a same-resolution domain translation task ($\times$1) that maps registered 3T acquisitions to their masked 7T counterparts, and a combined domain-translation and super-resolution task ($\times$2) in which spatially downsampled ($\times$0.5 trilinear zoom) 3T volumes are used as inputs and the full-resolution masked 7T volumes as targets. In both cases the ground truth is defined exclusively within the Circle of Willis mask $\mathcal{M}$, focusing supervision on the haemodynamically relevant intraluminal region. Bridging the real-world domain gap between 3T and 7T 4D Flow MRI requires paired, in-vivo acquisitions --- a prerequisite that has not previously been systematically addressed in the context of cerebrovascular 4D Flow.

Domain adaptation and cross-scanner generalisation are active areas of research in medical image analysis. Differences in MRI scanner platforms introduce systematic biases in image contrast, noise characteristics, and artefact patterns that degrade the performance of models trained on one domain when applied to another \cite{glocker2023domain,perone2019transductive}. In the specific context of multi-field-strength image synthesis, deep learning approaches have demonstrated the feasibility of learning 7T-like tissue contrast and spatial detail from 3T structural MRI \cite{qu2020synth7t,cui2024synth7t}, but no comparable work has addressed the fundamentally more complex case of quantitative, vector-valued 4D Flow MRI, where the network must predict physically consistent, three-directional velocity fields rather than scalar image intensities.

Physics-informed learning offers a principled framework to incorporate domain knowledge --- the governing equations of fluid mechanics and the physical constraints of MRI signal formation --- into the training objective, improving both the physical plausibility and generalisation of the learned haemodynamic representations \cite{kissas2020,ayed2023piml}. In the context of 4D Flow MRI, leveraging the divergence-free constraint of incompressible flow and the vessel-mask-aware structure of vascular velocity fields allows the network to focus its representational capacity on the haemodynamically relevant intraluminal regions, suppressing noise and artefacts in the extravascular background.

In this work, we present \textbf{NeuroFlow}, a supervised deep learning framework for joint 3T-to-7T domain translation and super-resolution of cerebrovascular 4D Flow MRI. The key contributions of this work are as follows:

\begin{enumerate}
    \item We acquire and process the first dataset of paired 3T/7T 4D Flow MRI studies in patients with cerebrovascular pathology, establishing a rigorous multi-stage registration and motion correction pipeline to achieve voxel-level correspondence between field strengths.

    \item We introduce a dual-branch vascular segmentation strategy combining a topology-aware nnU-Net ensemble \cite{isensee2021nnunet,yang2023topcow} with classical Frangi/Sato vesselness filtering \cite{frangi1998,sato1998} to generate high-quality Circle of Willis masks that guide mask-aware training.

    \item We adapt and extend the SR4DFlowNet architecture \cite{ferdian2020,elahmar2023} with a CBAM-based pre-upsampling attention module, yielding two task-specific models in the main manuscript: \textbf{CBAM-DNS} for same-resolution 3T$\to$7T domain translation ($\times$1) and \textbf{CBAM-SupRes} for joint domain translation and super-resolution ($\times$2). Extended masked-input, cascade, and broader architectural ablation results are reported in the supplementary material.

    \item We design a composite, mask-aware loss function that separately penalises velocity reconstruction error within and outside the vascular lumen, with an additional isotropic total variation penalty \cite{rudin1992tv} on extraluminal regions to suppress physiologically implausible background predictions.

    \item We demonstrate that the resulting framework achieves substantial reductions in relative speed error (RSE) and relative magnitude error (RME) within the Circle of Willis, representing a meaningful step towards virtual 7T quality cerebrovascular haemodynamics from routine 3T acquisitions.
\end{enumerate}

The remainder of this paper is organised as follows. Section~2 reviews related work on 4D Flow MRI processing, MRI super-resolution, and domain adaptation. Section~3 describes the methodology in detail, covering data acquisition, registration, network architecture, loss formulation, and training protocol. Section~4 presents quantitative and qualitative results across architectural variants and task configurations. Section~5 discusses the clinical implications, limitations, and directions for future work.


\section*{Methodology}

\subsection*{Study population and 4D Flow MRI acquisition}

Seven patients with cerebrovascular pathology were prospectively
enrolled and scanned on two MRI platforms within a single imaging
session: a clinical 3\,T Siemens system and an investigational 7\,T
Siemens system. For each subject, paired 4D Flow MRI studies were
acquired on both scanners to minimise inter-session physiological
variability. Each acquisition produced four co-registered volumetric
time series --- a magnitude image and three directional phase-contrast
velocity maps encoding velocity along the feet--head ($u$),
left--right ($v$), and anterior--posterior ($w$) axes
\cite{markl2012,dyverfeldt2015}. The 3\,T series covered between nine
and fourteen cardiac frames, while the 7\,T series spanned ten to
fifteen frames, all encoded with a common velocity-encoding sensitivity
of $\mathrm{VENC} = 0.9$\,\textnormal{m/s} across both field strengths
\cite{bock2010}. Raw data were converted from DICOM to NIfTI, preserving
the original phase representation and applying the LPS\,$\rightarrow$\,RAS
sign convention so that the three velocity components retained their
physical orientation in the scanner frame.

% FIGURE — Acquisition and pairing overview.
\begin{figure}[!htbp]
    \centering
    \includegraphics[width=0.8\textwidth]{Figures/figure1_acquisition_pairing_distill.pdf}
    \caption{Acquisition and pairing overview. Paired 3\,T and 7\,T 4D
    Flow MRI acquired within a single imaging session for each subject.}
    \label{fig:acquisition}
\end{figure}

\subsection*{Motion correction and multi-field-strength registration}

Cardiac-gated 4D Flow MRI is intrinsically sensitive to bulk motion
between cardiac phases. Intra-scan motion was corrected by estimating
rigid-body transformations $\{\mathbf{T}_t^{\text{rigid}}\}_{t=1}^{T}$
from the magnitude time series relative to a reference frame using
normalised mutual information \cite{maes1997}, and applying each
transformation consistently to the paired velocity volume,
\begin{equation}
    \mathbf{V}_t^{\text{corr}}
    = \mathbf{T}_t^{\text{rigid}} \circ \mathbf{V}_t^{\text{raw}},
    \qquad t = 1, \ldots, T,
\end{equation}
followed by a vector-aware reorientation of the velocity components. To
bring the 7\,T volumes into the geometric space of the 3\,T
acquisitions, a temporal maximum-intensity projection
$\mathrm{MIP}(\mathbf{x}) = \max_t M_t(\mathbf{x})$ was first computed
for each magnitude series to obtain a static, high-contrast angiographic
representation. A rigid alignment followed by an affine refinement,
both optimised with normalised mutual information, defined the
composite transformation $\mathbf{T}^{\text{affine}}_{7T \to 3T}$, which
was then applied to the full 4D velocity volume so that
\begin{equation}
    \mathbf{V}_{7T \to 3T}^{\text{reg}}(\mathbf{x}, t)
    = \mathbf{T}^{\text{affine}}_{7T \to 3T} \circ
      \mathbf{V}_{7T}^{\text{corr}}(\mathbf{x}, t),
\end{equation}
establishing voxel-level correspondence between the paired 3\,T and
7\,T data as required for supervised learning.

% FIGURE — Registration pipeline.
\begin{figure}[!htbp]
    \centering
    \includegraphics[width=\textwidth]{Figures/figure2_registration_simplified.pdf}
    \caption{Registration pipeline. Temporal MIPs of the paired 3\,T and
    7\,T magnitude before and after the rigid + affine inter-scan
    registration; the checkerboard overlay illustrates the resulting
    voxel-level alignment.}
    \label{fig:registration}
\end{figure}

\subsection*{Angiographic volume synthesis and vascular masking}

A voxel-level vascular probability map was constructed by combining
two complementary contrast sources. A 4D speed image was obtained from
the VENC-scaled velocity vector,
\begin{equation}
    S(\mathbf{x}, t) = \sqrt{
        \left(\tfrac{u(\mathbf{x},t)}{\mathrm{VENC}}\right)^2 +
        \left(\tfrac{v(\mathbf{x},t)}{\mathrm{VENC}}\right)^2 +
        \left(\tfrac{w(\mathbf{x},t)}{\mathrm{VENC}}\right)^2},
\end{equation}
and projected over the cardiac cycle as
$\bar{S}(\mathbf{x}) = \max_t S(\mathbf{x},t)$. The temporal MIP of the
magnitude and $\bar{S}$ were each min--max normalised and their
voxel-wise product
$\mathcal{A}(\mathbf{x}) = \hat{\mathrm{MIP}}(\mathbf{x}) \cdot
\hat{S}(\mathbf{x})$ defined a synthetic angiographic volume in which
voxels that are both bright in magnitude and fast in flow are reinforced
while static tissue and background noise are suppressed.

The Circle-of-Willis (CoW) mask used throughout this work was generated
on the 7\,T magnitude data, taking advantage of their substantially
higher signal-to-noise ratio. For every subject, an initial foreground
estimate was obtained by adaptive intensity thresholding of the 7\,T
maximum-intensity projection with connected-component cleanup, then
refined by an ensemble nnU-Net inference and finally corrected by an
expert reviewer on a slice-by-slice basis. The resulting manually
corrected mask defined the gold-standard CoW segmentation in the 7\,T
geometry; transferring it to the 3\,T frame via the registration above
yielded the binary mask $\mathcal{M}(\mathbf{x})$ used as the
supervision and evaluation region.

Two nnU-Net models with the residual-encoder \emph{ResEnc-M} backbone
\cite{isensee2021nnunet} were then trained on these corrected labels
with the skeleton-recall loss formulation introduced by the
TopCoW-2024 winning solution \cite{yang2023topcow,aydin2024topcow}.
The objective combined Dice, skeleton-recall, and cross-entropy terms,
\begin{equation}
    \mathcal{L}_{\text{seg}}
    = w_{D}\,\mathcal{L}_{\text{Dice}}
    + w_{S}\,\mathcal{L}_{\text{SkelRec}}
    + w_{C}\,\mathcal{L}_{\text{CE}},
    \qquad w_{D}=w_{S}=w_{C}=1,
\end{equation}
where the skeleton-recall term is computed on the fly from a
differentiable topological skeletonisation of the ground-truth mask.
Both 2D and 3D-full-resolution configurations were trained over five
cross-validation folds and combined as a 2D\,$+$\,3D ensemble at
inference. The first model operates on the temporal maximum projection
of the native 3\,T magnitude (mean cross-validated Dice~$0.83$), and a
second model was trained on the same projection after isotropic
$0.5{\times}$ downsampling (mean Dice~$0.76$), matching the input
geometry of the super-resolution task introduced below. The same
ensemble is used twice in this work: first, at training time, to
propagate the 7\,T-derived gold-standard mask through the 3\,T,
$0.5{\times}$\,3\,T, and 7\,T-in-3\,T volumes that form the supervised
pairs; and second, at evaluation time, to re-segment each predicted
4D Flow volume so that its CoW geometry can be compared against the
subject-specific 7\,T reference mask.

% FIGURE — Angiographic synthesis and vascular masking.
\begin{figure}[!htbp]
    \centering
    \includegraphics[width=\textwidth]{Figures/figure3_simplified_white_3d_mask.pdf}
    \caption{Angiographic synthesis and vascular masking.
(\textit{Top}) The 3T temporal magnitude maximum-intensity projection
is multiplied voxel-wise by the normalised speed projection to form
the angiographic volume $A(\mathbf{x})$.
A binary Circle-of-Willis mask $M$ was derived from the 7T reference
magnitude by intensity thresholding, a bootstrap nnU-Net ensemble,
and slice-by-slice manual correction, yielding the ground-truth labels
used throughout this work.
(\textit{Bottom}) Two nnU-Net ResEnc-M models were trained with the
compound Skeleton-Recall loss
$\mathcal{L} = \mathcal{L}_{\mathrm{Dice}} + \mathcal{L}_{\mathrm{SkelRec}} + \mathcal{L}_{\mathrm{CE}}$
on the corrected masks: one on native 3T magnitude projections
(mean cross-validated Dice~$0.83$) and one on
$0.5\times$-downsampled 3T projections
(Dice~$0.76$).
The resulting inference masks $\hat{M}$ define the fluid region for
the super-resolution loss and for all downstream evaluation metrics.}
    \label{fig:segmentation}
\end{figure}

\subsection*{Network architecture: CBAM-DNS and CBAM-SupRes}

The core predictive backbone is an adapted SR4DFlowNet
\cite{ferdian2020,elahmar2023}, a 3D fully convolutional residual
network inspired by enhanced residual architectures for super-resolution
\cite{lim2017edsr}, augmented with a Convolutional Block Attention
Module (CBAM) \cite{woo2018cbam} inserted between the low-resolution
residual stack and the upsampling stage. We instantiate the same
backbone in two task-specific forms used throughout the results:
\textbf{CBAM-DNS} for same-resolution 3\,T$\to$7\,T domain translation
($r=1$) and \textbf{CBAM-SupRes} for joint domain translation and
twofold super-resolution ($r=2$).

The model accepts six scalar input channels per voxel --- the three
VENC-normalised velocity components $(u, v, w)$ and their
magnitude-weighted counterparts $(u_m, v_m, w_m)$ --- and internally
derives three auxiliary fields that summarise the local kinematics and
phase-contrast signal:
\begin{equation}
    S = \sqrt{u^{2}+v^{2}+w^{2}},
    \qquad
    M = \sqrt{u_m^{2}+v_m^{2}+w_m^{2}},
    \qquad
    \mathrm{PCMR} = M \cdot S.
\end{equation}
Velocity and physical channels are processed in two dedicated
convolutional streams of two $3{\times}3{\times}3$ layers with 64
feature maps, symmetric padding, and ReLU activation: a phase stream
operating on $(u,v,w)$ to extract flow-kinematic features, and a
phase-contrast stream operating on $(\mathrm{PCMR}, M, S)$ to extract
signal-intensity features. Their feature maps are concatenated and
fused by a $1{\times}1{\times}1$ projection followed by a
$3{\times}3{\times}3$ refinement convolution.

The fused representation is propagated through $N_{LR}=8$ pre-activation
residual blocks at the input resolution \cite{he2016resnet},
\begin{equation}
    \mathbf{F}^{(l+1)} = \mathbf{F}^{(l)}
        + f\!\left(\mathbf{F}^{(l)};\, \Theta^{(l)}\right),
\end{equation}
after which the CBAM module sequentially applies channel and spatial
attention to re-weight feature responses before upsampling. Channel
attention is obtained from global average- and max-pooled descriptors
processed by a shared two-layer multilayer perceptron, while spatial
attention is computed from the channel-wise pooled feature maps and a
$7{\times}7{\times}7$ convolution; the two attention maps are applied
multiplicatively in sequence. Spatial upsampling is then performed by
trilinear interpolation with scale factor $r$, followed by $N_{HR}=4$
high-resolution residual blocks operating at the target resolution. The
network terminates in three independent two-layer convolutional heads
that emit the reconstructed velocity components
$(\hat{u}, \hat{v}, \hat{w})$, together with an additional head that
predicts the reconstructed magnitude $\hat{M}$. Within this common
backbone, CBAM-DNS uses $r=1$ and trains under the same-resolution
domain-translation setting, whereas CBAM-SupRes uses $r=2$ and absorbs
both the cross-field-strength translation and the spatial upsampling
into a single forward pass.

% FIGURE — Network architecture diagram.
\begin{figure}[!htbp]
    \centering
    \includegraphics[width=\textwidth]{Figures/Figure_4_CBAM_DNS_SupRes.pdf}
    \caption{Network architecture underlying CBAM-DNS ($r=1$) and
    CBAM-SupRes ($r=2$). The phase and phase-contrast streams are fused,
    refined by $N_{LR}=8$ residual blocks, re-weighted by a CBAM
    channel-and-spatial attention module, upsampled by trilinear
    interpolation with factor $r$, and refined again by $N_{HR}=4$
    high-resolution residual blocks before three velocity heads and an
    optional magnitude head emit the reconstructed volume.}
    \label{fig:architecture}
\end{figure}

\subsection*{Loss function and evaluation metrics}

Training minimises a composite, mask-aware loss that jointly penalises
intraluminal velocity error, magnitude error, and texture irregularity
in the extraluminal background,
\begin{equation}
    \mathcal{L} = \mathcal{L}_{\text{vel}}
                + \lambda_{M}\,\mathcal{L}_{\text{mag}}
                % + \lambda_{\text{TV}}\,\mathcal{L}_{\text{TV}}.
\end{equation}
For each predicted velocity component
$\hat{q}\in\{\hat{u},\hat{v},\hat{w}\}$ and its VENC-normalised target
$q^{*}=q/\mathrm{VENC}$, the per-voxel squared error
$\epsilon_q(\mathbf{x})=(\hat{q}(\mathbf{x})-q^{*}(\mathbf{x}))^{2}$ is
split into a masked fluid term and a complementary background term,
\begin{equation}
    \mathcal{L}^{q}_{\text{fluid}}
    = \frac{\sum_{\mathbf{x}} \epsilon_q(\mathbf{x})\,\mathcal{M}(\mathbf{x})}
           {\sum_{\mathbf{x}} \mathcal{M}(\mathbf{x}) + \varepsilon},
    \qquad
    \mathcal{L}^{q}_{\text{bg}}
    = \frac{\sum_{\mathbf{x}} \epsilon_q(\mathbf{x})\,(1-\mathcal{M}(\mathbf{x}))}
           {\sum_{\mathbf{x}} (1-\mathcal{M}(\mathbf{x})) + \varepsilon},
\end{equation}
and aggregated as
$\mathcal{L}_{\text{vel}}
 = \sum_{q\in\{u,v,w\}}
   (\mathcal{L}^{q}_{\text{fluid}}+\lambda_{\text{bg}}\,\mathcal{L}^{q}_{\text{bg}})$
with background weight $\lambda_{\text{bg}}=0.1$, reflecting the
clinical priority of accurate intraluminal velocity quantification.
$\mathcal{L}_{\text{mag}}$ is the analogous masked mean-squared error on
the per-frame min--max-normalised magnitude with magnitude weight
$\lambda_{M}=1$.
% To suppress spurious high-frequency texture in
% extraluminal regions, an isotropic total-variation penalty
% \cite{rudin1992tv} is applied exclusively to voxel pairs that lie
% outside the mask along each Cartesian axis,
% \begin{equation}
%     \mathcal{L}_{\text{TV}}
%     = \sum_{q}\sum_{d\in\{x,y,z\}}
%       \frac{\sum_{\mathbf{x}} \lvert
%       \hat{q}(\mathbf{x}+\mathbf{e}_d)-\hat{q}(\mathbf{x}) \rvert\,
%       \Omega(\mathbf{x},d)}
%            {\sum_{\mathbf{x}} \Omega(\mathbf{x},d) + \varepsilon},
% \end{equation}
% where
% $\Omega(\mathbf{x},d) =
%  (1-\mathcal{M}(\mathbf{x}))\,(1-\mathcal{M}(\mathbf{x}+\mathbf{e}_d))$
% is the joint outside indicator and
% $\lambda_{\text{TV}}=10^{-5}$ provides a gentle smoothness prior
% without over-regularising the extraluminal signal.

Image-reconstruction quality was assessed by VENC-normalised
mean-squared error (MSE) of the predicted velocity field. The
intraluminal velocity error was computed inside the CoW mask both
as a total
\begin{equation}
    \mathrm{MSE}_{\text{vel}}
    = \frac{\sum_{\mathbf{x}}
            \sum_{q\in\{u,v,w\}}\!\!
            (\hat{q}(\mathbf{x})-q^{*}(\mathbf{x}))^{2}\,
            \mathcal{M}(\mathbf{x})}
           {\sum_{\mathbf{x}}\mathcal{M}(\mathbf{x}) + \varepsilon}
\end{equation}
and per component
\begin{equation}
    \mathrm{MSE}_{q}
    = \frac{\sum_{\mathbf{x}}
            (\hat{q}(\mathbf{x})-q^{*}(\mathbf{x}))^{2}\,
            \mathcal{M}(\mathbf{x})}
           {\sum_{\mathbf{x}}\mathcal{M}(\mathbf{x}) + \varepsilon},
    \qquad q \in \{u,v,w\},
\end{equation}
and an analogous extraluminal quantity --- the background MSE
$\mathrm{MSE}_{\text{bg}}$ --- was obtained by substituting
$\mathcal{M}$ with $1-\mathcal{M}$ in the numerator and denominator.
Both metrics are reported in units of $10^{-3}$ in the
VENC-normalised range.

Vascular-geometry fidelity was evaluated by re-segmenting each
predicted 4D Flow volume with the same ResEnc-M ensemble described
above and comparing the resulting binary CoW mask $P$ to the
subject-specific 7\,T reference mask $R$ in the 3\,T geometry. Volumetric
overlap was quantified by the Dice score,
\begin{equation}
    \mathrm{Dice}(P,R)
    = \frac{2\lvert P \cap R \rvert}{\lvert P \rvert + \lvert R \rvert}.
\end{equation}
Boundary agreement was quantified in physical units from
triangulated surfaces $S_P$ and $S_R$ extracted by marching cubes.
Letting
$d(p,S_R) = \min_{r \in S_R} \lVert p - r \rVert_{2}$ denote the
directed point-to-surface distance, the symmetric mean surface
distance is
\begin{equation}
    \mathrm{sMSD}(P,R)
    = \tfrac{1}{2}\!\left(
        \frac{1}{|S_P|}\sum_{p\in S_P} d(p,S_R)
        + \frac{1}{|S_R|}\sum_{r\in S_R} d(r,S_P)
      \right),
\end{equation}
and the worst-case boundary mismatch is captured by the Hausdorff
distance,
\begin{equation}
    \mathrm{HD}(P,R)
    = \max\!\left(
        \max_{p\in S_P} d(p,S_R),\;
        \max_{r\in S_R} d(r,S_P)
      \right).
\end{equation}
Dice is reported as higher-is-better in $[0,1]$; sMSD and HD are
reported in millimetres and lower-is-better. Together, these metrics
quantify both the volumetric overlap and the boundary fidelity of the
reconstructed vascular geometry that the velocity field implies.


\subsection*{Data preparation, training, and inference}

Prior to network input, all velocity components were normalised by the
global VENC,
\begin{equation}
    \tilde{q} = q/\mathrm{VENC}_{\text{global}},
    \qquad
    \mathrm{VENC}_{\text{global}}
    = \max(\mathrm{VENC}_u,\mathrm{VENC}_v,\mathrm{VENC}_w),
\end{equation}
yielding values approximately in $[-1,1]$, and magnitude images were
rescaled to $[0,1]$ by per-frame min--max scaling
\cite{cardoso2022monai}. In both task configurations the reconstruction
target was the registered 7\,T velocity and magnitude volume after
masking by $\mathcal{M}$,
$\mathbf{V}^{*}(\mathbf{x},t)=\mathbf{V}^{\text{reg}}_{7T}(\mathbf{x},t)\,
\mathcal{M}(\mathbf{x})$, confining supervision to the haemodynamically
relevant intraluminal region. CBAM-DNS was trained at native 3\,T
resolution ($r=1$), with $\mathbf{V}^{\text{reg}}_{3T}$ as input and
$\mathbf{V}^{*}$ as target, so that the network simultaneously bridges
the SNR gap between the two field strengths, corrects residual contrast
differences arising from distinct hardware and reconstruction chains,
and recovers finer vascular detail visible only at 7\,T. CBAM-SupRes
was trained with a spatially downsampled input
$\mathbf{V}^{\mathrm{LR}}_{3T}(\mathbf{x},t)
 = \mathcal{Z}_{0.5}[\mathbf{V}^{\text{reg}}_{3T}(\mathbf{x},t)]$,
obtained by isotropic trilinear downsampling, and the same masked
7\,T target, combining $r=2$ upsampling and domain translation in a
single end-to-end pass.

Each task was trained in two complementary input strategies that are
reported side by side in the results. In the \emph{synthetic-noise}
strategy, the network sees the original 3\,T (or downsampled 3\,T)
volume with synthetic non-vascular noise injected into the extraluminal
region during training, and operates without a mask at inference time.
In the \emph{masked-input} strategy, the predicted CoW segmentation is
used to zero the extravascular signal of the 3\,T input prior to the
network for both training and inference; no noise augmentation is then
applied. The synthetic-noise strategy targets mask-free deployment and
emphasises robustness to background patterns, whereas the masked-input
strategy emphasises geometric fidelity by removing the extraluminal
domain shift altogether.

Two augmentation strategies were applied online. Vector-aware geometric
augmentation drew random 90\textdegree{} rotations from the planes
$\{xy,xz,yz\}$ with multiplicities $\{1,2,3\}$ at probability $0.5$,
and the three velocity components were permuted and sign-corrected as
a proper vector field while the mask and magnitude were rotated as
scalar fields. Synthetic non-vascular noise augmentation (used in the
synthetic-noise strategy only) injected per-batch amplitudes drawn
from one of several statistical distributions (zero-mean normal,
Student-$t$, Laplace, skew-normal, or generalised normal) into the
extraluminal region of the input,
\begin{equation}
    \tilde{\mathbf{V}}^{\mathrm{LR}}(\mathbf{x})
    = \mathbf{V}^{\mathrm{LR}}(\mathbf{x})
      + \eta(\mathbf{x})\,(1-\mathcal{M}(\mathbf{x})),
\end{equation}
encouraging mask-focused learning and preventing the network from
overfitting to specific noise patterns in the background.

Training was conducted on three-dimensional patches of size $16^{3}$
voxels in the input space, which corresponded to $32^{3}$ voxels in
the output space for the $\times 2$ configuration
\cite{peiffer2024patch}. For each case, 64 patches were drawn per epoch
with one cardiac frame selected uniformly at random per sample; 16
deterministic centre patches were used for validation loss monitoring,
and an optional minimum-mask-coverage criterion ensured that each
sampled patch contained a sufficient proportion of vascular voxels.
Optimisation used Adam \cite{kingma2015adam} with an initial learning
rate of $2{\times}10^{-4}$ and weight decay $5{\times}10^{-7}$. A
ReduceLROnPlateau scheduler monitored the validation loss with a
patience of ten epochs and halved the learning rate (floor $10^{-6}$),
and an early-stopping criterion with patience twenty epochs,
complemented by an overfitting detector that terminated training when
the training loss kept decreasing while the validation loss increased
for fifteen consecutive epochs, prevented over-fitting. Models were
trained with mini-batches of four samples; the checkpoint attaining the
lowest validation loss was retained for inference.

All results are reported under \emph{leave-one-out cross-validation}
(LOOCV) at the patient level. With $n=7$ subjects, seven independent
folds were trained: in each fold, one subject was held out as the test
case and the remaining six formed the training set; no subject appeared
in both subsets. The metrics reported in the results are mean $\pm$
standard deviation across the seven held-out subjects.

At inference time, full 4D volumes were reconstructed with an
overlap-aware sliding-window strategy
\cite{cardoso2022monai,ronneberger2015unet}. For each time frame the
volume was partitioned into overlapping $16^{3}$ patches with $25\%$
overlap, and predicted patches were recombined by Gaussian-weighted
blending at overlapping boundaries,
\begin{equation}
    \hat{\mathbf{V}}(\mathbf{x})
    = \frac{\sum_{k} w_{k}(\mathbf{x})\,\hat{\mathbf{V}}_{k}(\mathbf{x})}
           {\sum_{k} w_{k}(\mathbf{x})},
\end{equation}
where $w_k(\mathbf{x})$ is a Gaussian window centred on patch $k$,
eliminating patch-edge discontinuities. The reconstructed volumes were
masked by $\mathcal{M}$ to constrain the output to the CoW region for
the intraluminal metrics, and the predicted velocity components were
finally rescaled to physical units via
$\hat{q}_{\text{phys}} = \hat{q}\cdot\mathrm{VENC}$. For the geometry
evaluation, the reconstructed magnitude (and, when unavailable, the
synthesised angiographic volume of Figure~\ref{fig:segmentation}) was passed
through the ResEnc-M segmentation ensemble to produce the predicted
binary CoW mask $P$ that was subsequently compared against the
subject-specific 7\,T reference $R$ to obtain Dice, sMSD, and HD.


% =====================================================================
%  END OF METHODOLOGY SECTION
% =====================================================================

% ============================================================
%  Section 4 --- Results
% ============================================================
\section*{Results}

\subsection*{4.1 Experimental Overview}

The main manuscript evaluates two input strategies across two task modes. In the
\textbf{synthetic-noise} setting (mask-free deployment), the 3T input is fed directly to the
network after synthetic noise augmentation. In the \textbf{masked-input} setting, a
predicted CoW segmentation from nnUNet is used to zero extravascular signal in the 3T volume
prior to reconstruction; no noise augmentation is applied. Six configurations are retained
in the primary results (Table~\ref{tab:experiment_map}).

All results are reported as Leave-One-Out Cross-Validation (LOOCV) means $\pm$ standard
deviation across seven subjects. Geometry metrics are computed by re-segmenting each predicted
4D Flow volume with the ToP-CoW pipeline and comparing against the subject-specific 7T
reference mask. Bold indicates the best mean value within each task group across both input
strategies.

\begin{table}[!htbp]
  \centering
  \caption{Configurations retained in the main manuscript.}
  \label{tab:experiment_map}
  \small
  \begin{tabular}{clll}
    \toprule
    ID & Architecture & Input strategy & Noise aug. \\
    \midrule
    \multicolumn{4}{l}{\textit{Denoising ($\times 1$)}} \\
    \midrule
    A & Original  & Synthetic noise & Yes \\
    B & CBAM-DNS  & Synthetic noise & Yes \\
    E & CBAM-DNS  & Masked input    & No  \\
    \midrule
    \multicolumn{4}{l}{\textit{Super-Resolution ($\times 2$)}} \\
    \midrule
    C & Original    & Synthetic noise & Yes \\
    D & CBAM-SupRes & Synthetic noise & Yes \\
    F & CBAM-SupRes & Masked input    & No  \\
    G & 4DFlowNet   & Original input  & No  \\
    \bottomrule
  \end{tabular}
\end{table}

% -------------------------------------------------------
\subsection*{4.2 Velocity Reconstruction}
% -------------------------------------------------------
\label{sec:velocity}

\begin{table}[!htbp]
  \centering
  \caption{Intraluminal velocity reconstruction metrics (LOOCV, $n=7$; mean\,$\pm$\,SD).
  VENC-normalised MSE ($\times 10^{-3}$; lower is better). Bold: best mean within each
  task group.}
  \label{tab:velocity}
  \setlength{\tabcolsep}{4pt}
  \small
  \begin{tabular}{clll rrrr}
    \toprule
    \multirow{2}{*}{ID} & \multirow{2}{*}{Architecture} & \multirow{2}{*}{Input} &
      \multicolumn{4}{c}{MSE $\downarrow$ ($\times 10^{-3}$)} \\
    \cmidrule(lr){4-7}
    & & & Total & $u$ & $v$ & $w$ \\
    \midrule
    \multicolumn{7}{l}{\textit{Denoising ($\times 1$) --- Synthetic noise}} \\
    \midrule
    A & Original  & Synth. & $35.9 \pm 11.5$ & $11.3 \pm 3.8$ & $14.2 \pm 4.6$ & $10.5 \pm 3.6$ \\
    B & CBAM-DNS  & Synth. & $35.7 \pm 12.6$ & $11.6 \pm 3.8$ & $13.5 \pm 5.3$ & $10.6 \pm 3.9$ \\
    \multicolumn{7}{l}{\textit{Denoising ($\times 1$) --- Masked input}} \\
    \midrule
    E & CBAM-DNS  & Masked & $\mathbf{34.5 \pm 11.6}$ & $\mathbf{11.2 \pm 3.7}$ & $\mathbf{13.0 \pm 4.8}$ & $\mathbf{10.2 \pm 3.6}$ \\
    \midrule
    \multicolumn{7}{l}{\textit{Super-Resolution ($\times 2$) --- Synthetic noise}} \\
    \midrule
    C & Original    & Synth. & $41.5 \pm 15.1$ & $13.3 \pm 4.4$ & $15.4 \pm 6.2$ & $12.7 \pm 5.0$ \\
    D & CBAM-SupRes & Synth. & $\mathbf{40.1 \pm 14.3}$ & $\mathbf{12.9 \pm 4.4}$ & $\mathbf{14.8 \pm 5.7}$ & $12.5 \pm 4.7$ \\
    G & 4DFlowNet   & Orig.  & $62.1 \pm 21.4$ & $20.7 \pm 6.9$ & $22.4 \pm 8.2$ & $19.0 \pm 6.7$ \\
    \multicolumn{7}{l}{\textit{Super-Resolution ($\times 2$) --- Masked input}} \\
    \midrule
    F & CBAM-SupRes & Masked & $40.8 \pm 13.9$ & $13.0 \pm 4.2$ & $15.5 \pm 5.7$ & $\mathbf{12.3 \pm 4.6}$ \\
    \bottomrule
  \end{tabular}
\end{table}

In the denoising task, \textbf{CBAM-DNS} with masked input achieves the lowest total
intraluminal MSE ($34.5 \pm 11.6\,{\times}10^{-3}$), a 3.4\% improvement over synthetic-noise
CBAM-DNS and 3.9\% over the Original baseline. In the super-resolution task,
\textbf{CBAM-SupRes} trained with synthetic-noise augmentation retains the lowest total MSE
($40.1 \pm 14.3\,{\times}10^{-3}$) and the best $u$ and $v$ components, while the masked-input
variant leads only on the $w$ component. The original \textbf{4DFlowNet} baseline, although
trained specifically for $\times 2$ super-resolution, yields substantially higher error
($62.1 \pm 21.4\,{\times}10^{-3}$), indicating weaker generalisation to this 3T$\rightarrow$7T
brain angiography setting than the NeuroFlow variants.

% -------------------------------------------------------
\subsection*{4.3 Vascular Geometry Fidelity}
% -------------------------------------------------------

\begin{table}[!htbp]
  \centering
  \caption{Vascular geometry metrics (LOOCV, $n=7$; mean\,$\pm$\,SD). Dice ($\uparrow$);
  sMSD in mm ($\downarrow$); HD in mm ($\downarrow$). Bold: best mean within each task group.}
  \label{tab:geometry}
  \setlength{\tabcolsep}{4.5pt}
  \small
  \begin{tabular}{clll rrr}
    \toprule
    ID & Architecture & Input & Dice $\uparrow$ & sMSD $\downarrow$ (mm) & HD $\downarrow$ (mm) \\
    \midrule
    \multicolumn{6}{l}{\textit{Denoising ($\times 1$) --- Synthetic noise}} \\
    \midrule
    A & Original  & Synth. & $0.826 \pm 0.039$ & $0.376 \pm 0.088$ & $14.67 \pm 3.75$ \\
    B & CBAM-DNS  & Synth. & $0.831 \pm 0.032$ & $0.359 \pm 0.074$ & $\mathbf{13.67 \pm 5.28}$ \\
    \multicolumn{6}{l}{\textit{Denoising ($\times 1$) --- Masked input}} \\
    \midrule
    E & CBAM-DNS  & Masked & $\mathbf{0.839 \pm 0.022}$ & $\mathbf{0.300 \pm 0.037}$ & $14.12 \pm 4.85$ \\
    \midrule
    \multicolumn{6}{l}{\textit{Super-Resolution ($\times 2$) --- Synthetic noise}} \\
    \midrule
    C & Original    & Synth. & $0.778 \pm 0.037$ & $0.434 \pm 0.063$ & $15.28 \pm 3.73$ \\
    D & CBAM-SupRes & Synth. & $0.781 \pm 0.041$ & $0.417 \pm 0.077$ & $14.53 \pm 3.92$ \\
    \multicolumn{6}{l}{\textit{Super-Resolution ($\times 2$) --- Masked input}} \\
    \midrule
    F & CBAM-SupRes & Masked & $\mathbf{0.786 \pm 0.033}$ & $\mathbf{0.370 \pm 0.049}$ & $\mathbf{14.01 \pm 4.95}$ \\
    \bottomrule
  \end{tabular}
\end{table}

The masked-input variant improves Dice and sMSD in both tasks. In denoising, \textbf{CBAM-DNS}
(Masked) achieves the best Dice ($0.839 \pm 0.022$) and sMSD ($0.300 \pm 0.037$\,mm), with
the sMSD standard deviation roughly halved relative to synthetic-noise CBAM-DNS ($\pm 0.037$
vs.\ $\pm 0.074$), reflecting more consistent vessel boundary recovery across subjects.
HD is best for synthetic-noise CBAM-DNS ($13.67$\,mm), indicating that masking reduces
average surface error but not necessarily the worst-case outlier distance. In the
super-resolution task, \textbf{CBAM-SupRes} (Masked) leads on all three geometry metrics.
Geometry metrics for 4DFlowNet were not included because the corresponding segmentation-based
evaluation was not available for this baseline at the time of analysis.

% -------------------------------------------------------
\subsection*{4.4 Background Suppression}
% -------------------------------------------------------

\begin{table}[!htbp]
  \centering
  \caption{Background velocity MSE outside the CoW mask (LOOCV, $n=7$; mean\,$\pm$\,SD).
  VENC-normalised MSE ($\times 10^{-3}$; lower is better). Bold: best mean within each
  task group.}
  \label{tab:background}
  \setlength{\tabcolsep}{4pt}
  \small
  \begin{tabular}{clll rrrr}
    \toprule
    ID & Architecture & Input &
      $\text{MSE}_{\text{bg}}$ $\downarrow$ &
      $\text{MSE}_{u}$ $\downarrow$ &
      $\text{MSE}_{v}$ $\downarrow$ &
      $\text{MSE}_{w}$ $\downarrow$ \\
    & & & \multicolumn{4}{c}{($\times 10^{-3}$, mean\,$\pm$\,SD)} \\
    \midrule
    \multicolumn{7}{l}{\textit{Denoising ($\times 1$) --- Synthetic noise}} \\
    \midrule
    A & Original  & Synth. & $0.622 \pm 0.191$ & $0.234 \pm 0.096$ & $0.234 \pm 0.076$ & $0.154 \pm 0.046$ \\
    B & CBAM-DNS  & Synth. & $0.672 \pm 0.248$ & $0.262 \pm 0.105$ & $0.257 \pm 0.100$ & $0.154 \pm 0.052$ \\
    \multicolumn{7}{l}{\textit{Denoising ($\times 1$) --- Masked input}} \\
    \midrule
    E & CBAM-DNS  & Masked & $\mathbf{0.338 \pm 0.099}$ & $\mathbf{0.126 \pm 0.034}$ & $\mathbf{0.130 \pm 0.047}$ & $\mathbf{0.081 \pm 0.022}$ \\
    \midrule
    \multicolumn{7}{l}{\textit{Super-Resolution ($\times 2$) --- Synthetic noise}} \\
    \midrule
    C & Original    & Synth. & $0.916 \pm 0.664$ & $0.361 \pm 0.303$ & $0.356 \pm 0.256$ & $0.199 \pm 0.109$ \\
    D & CBAM-SupRes & Synth. & $0.588 \pm 0.152$ & $0.202 \pm 0.046$ & $0.243 \pm 0.084$ & $0.144 \pm 0.034$ \\
    G & 4DFlowNet   & Orig.  & $65.462 \pm 24.410$ & $21.990 \pm 8.167$ & $21.832 \pm 8.160$ & $21.640 \pm 8.085$ \\
    \multicolumn{7}{l}{\textit{Super-Resolution ($\times 2$) --- Masked input}} \\
    \midrule
    F & CBAM-SupRes & Masked & $\mathbf{0.346 \pm 0.113}$ & $\mathbf{0.139 \pm 0.040}$ & $\mathbf{0.122 \pm 0.053}$ & $\mathbf{0.084 \pm 0.025}$ \\
    \bottomrule
  \end{tabular}
\end{table}

The masked-input variant achieves the lowest background MSE in all configurations by a large
margin. In the denoising setting, $\text{MSE}_{\text{bg}}$ drops to
$0.338\,{\times}10^{-3}$ ($-45.7\%$ vs.\ Original, $-49.7\%$ vs.\ synthetic CBAM-DNS). In
super-resolution, $0.346\,{\times}10^{-3}$ ($-41.2\%$ vs.\ CBAM-SupRes, $-62.2\%$ vs.\
Original). By contrast, the 4DFlowNet baseline shows markedly elevated background error
($65.462 \pm 24.410\,{\times}10^{-3}$), consistent with weaker suppression of non-vascular
signal under this domain shift. Consequently, background metrics for masked-input models
should be interpreted as a lower bound rather than a direct comparison against unconstrained
synthetic-noise models.

% -------------------------------------------------------
\subsection*{4.5 Consolidated Comparison}
% -------------------------------------------------------

\begin{table}[!htbp]
  \centering
  \caption{Consolidated LOOCV results ($n=7$; mean\,$\pm$\,SD). $\text{MSE}_\text{vel}$ and
  $\text{MSE}_\text{bg}$ in units of $10^{-3}$ (VENC-normalised); sMSD and HD in mm.
  Bold: best mean within each task group.}
  \label{tab:summary}
  \setlength{\tabcolsep}{3pt}
  \small
  \begin{tabular}{clll rrrrr}
    \toprule
    ID & Architecture & Input &
      \makecell{$\text{MSE}_\text{vel}$ $\downarrow$\\(${\times}10^{-3}$)} &
      Dice $\uparrow$ &
      \makecell{sMSD $\downarrow$\\(mm)} &
      \makecell{HD $\downarrow$\\(mm)} &
      \makecell{$\text{MSE}_\text{bg}$ $\downarrow$\\(${\times}10^{-3}$)} \\
    \midrule
    \multicolumn{8}{l}{\textit{Denoising ($\times 1$) --- Synthetic noise}} \\
    \midrule
    A & Original  & Synth. & $35.9 \pm 11.5$ & $0.826 \pm 0.039$ & $0.376 \pm 0.088$ & $14.67 \pm 3.75$ & $0.622 \pm 0.191$ \\
    B & CBAM-DNS  & Synth. & $35.7 \pm 12.6$ & $0.831 \pm 0.032$ & $0.359 \pm 0.074$ & $\mathbf{13.67 \pm 5.28}$ & $0.672 \pm 0.248$ \\
    \multicolumn{8}{l}{\textit{Denoising ($\times 1$) --- Masked input}} \\
    \midrule
    E & CBAM-DNS  & Masked & $\mathbf{34.5 \pm 11.6}$ & $\mathbf{0.839 \pm 0.022}$ & $\mathbf{0.300 \pm 0.037}$ & $14.12 \pm 4.85$ & $\mathbf{0.338 \pm 0.099}$ \\
    \midrule
    \multicolumn{8}{l}{\textit{Super-Resolution ($\times 2$) --- Synthetic noise}} \\
    \midrule
    C & Original    & Synth. & $41.5 \pm 15.1$ & $0.778 \pm 0.037$ & $0.434 \pm 0.063$ & $15.28 \pm 3.73$ & $0.916 \pm 0.664$ \\
    D & CBAM-SupRes & Synth. & $\mathbf{40.1 \pm 14.3}$ & $0.781 \pm 0.041$ & $0.417 \pm 0.077$ & $14.53 \pm 3.92$ & $0.588 \pm 0.152$ \\
    G & 4DFlowNet   & Orig.  & $62.1 \pm 21.4$ & N/A & N/A & N/A & $65.462 \pm 24.410$ \\
    \multicolumn{8}{l}{\textit{Super-Resolution ($\times 2$) --- Masked input}} \\
    \midrule
    F & CBAM-SupRes & Masked & $40.8 \pm 13.9$ & $\mathbf{0.786 \pm 0.033}$ & $\mathbf{0.370 \pm 0.049}$ & $\mathbf{14.01 \pm 4.95}$ & $\mathbf{0.346 \pm 0.113}$ \\
    \bottomrule
  \end{tabular}
\end{table}

Four principal conclusions emerge from the LOOCV results.

\textbf{CBAM-DNS and CBAM-SupRes with synthetic-noise training generalise consistently across
all seven subjects without requiring a segmentation mask at inference time.} In the
super-resolution task, synthetic-noise CBAM-SupRes retains the best velocity MSE
($40.1\,{\times}10^{-3}$), demonstrating that noise augmentation is the dominant factor for
intraluminal velocity accuracy when spatial upsampling is required.

\textbf{The masked-input variant improves geometry and background suppression in both tasks.}
CBAM-DNS (Masked) reduces sMSD by 16.4\% and its standard deviation by 50\% relative to
synthetic-noise CBAM-DNS, indicating more consistent vessel boundary recovery. The best
Hausdorff distance in denoising remains with synthetic-noise CBAM-DNS ($13.67$\,mm), as
masking reduces average surface error but not worst-case outlier distances.

\textbf{Background suppression is substantially lower for masked-input models by construction,}
and should not be directly compared to synthetic-noise models. Masking the 3T input prevents
extravascular signal from propagating through the network, making the background metric an
expected lower bound rather than an architectural advantage.

\textbf{The original 4DFlowNet baseline underperforms the NeuroFlow super-resolution models
on this task.} Although 4DFlowNet was originally designed for $\times 2$ super-resolution,
it yields higher intraluminal velocity error ($62.1 \pm 21.4\,{\times}10^{-3}$) and much
larger background error ($65.462 \pm 24.410\,{\times}10^{-3}$) than both the Original and
CBAM-SupRes NeuroFlow variants, suggesting limited robustness to the present 3T$\rightarrow$7T
brain angiography domain shift.

\textbf{Model selection depends on deployment context.} When a reliable CoW segmentation is
available at inference time, CBAM-DNS/SupRes (Masked) offers the best geometry and surface
fidelity. When operating mask-free, synthetic-noise CBAM-SupRes achieves the best velocity
reconstruction for the $\times 2$ task, and CBAM-DNS the best for native-resolution
denoising. Extended ablations and cascade results are provided in the supplementary material.

\bibliographystyle{ieeetr}
\bibliography{references}

\end{document}
